% RESEARCH_DRAFT_SHA256: 7f24b5c87ac10594adc65298a4d84d040ce2c52f92659c310820d1828b3c0f2e
% Layout and source-reference wrappers expanded for the web renderer.
\documentclass[11pt]{article}
\usepackage[margin=25mm]{geometry}
\usepackage{fontspec,unicode-math}
\setmainfont{Latin Modern Roman}
\setmathfont{Latin Modern Math}
\usepackage{amsmath,amssymb,mathtools,xurl,hyperref}
\hypersetup{colorlinks=true,urlcolor=blue,linkcolor=blue}
\setcounter{secnumdepth}{0}
\setlength{\parindent}{0pt}
\setlength{\parskip}{5pt}
\setlength{\emergencystretch}{3em}
\newcommand{\sref}[2]{\hyperlink{source:#1:#2}{[S#2]}}
\newcommand{\cataloguescope}[1]{\paragraph{Source record.}#1}
\newcommand{\floor}[1]{\left\lfloor#1\right\rfloor}
\begin{document}
\section*{Research-attempt provenance}
This standalone research attempt retains the complete original definition body below,
including its historical source and preparation notes. Those notes describe the
earlier catalogue preparation; the new attempt and its verification are identified
separately after that body. No original ID, catalogue status, or source statement is changed.

\section{Localization and Fourier coordinates beyond polynomial models}\label{statement}

\subsection{Definitions and mathematical statement}

\paragraph{Shared notation.}Standard number systems are $\mathbb N,\mathbb Z,\mathbb Q,\mathbb R,\mathbb C$. Specialized notation and index ranges follow the cited source and the definitions below.


A polynomial ring uses finite sums of monomials. Its rational function field consists of quotients of polynomials with nonzero denominators. Localization allows selected denominators to be inverted. Fourier analysis represents functions using characters or frequency components; the relevant algebra depends on the underlying group and function space.
\paragraph{Record scope and qualifications.}
The catalogue title is retained as metadata. Named certificates, contracts, bounds or frameworks in that title are not asserted as results of the original source extract.
No alternate ring, statistical model, or desired property is fixed. The question is a methodological extension agenda rather than a specified localization theorem. The supplied catalogue marks this record partially solved; its present research status was not independently checked.
The source text below is rendered literally where necessary; damaged formulas and composite extracts are not silently repaired into a new mathematical claim.
\paragraph{Statement recorded in the supplied source.}
Question 49 (Wynn). Why polynomials? What can be said about algebraic methods when applied to other rings, like rings of rational functions, or rings where Fourier analysis can be used? 11 \href{https://aimath.org/WWN/compalgstat/compalgstat.pdf}{[S2]}
\paragraph{Source status.}
The supplied catalogue labels this source record partially solved; that metadata was checked against the supplied snapshot on 2026-10-11. The label is preserved for provenance and is not an independent verification of present open status or mathematical correctness.
\subsection{Short English statement}
The source proposes extending algebraic statistical methods beyond polynomial coordinates.
\subsection{Sources}
\begin{itemize}
\item[S1] Ulam AI and the UnsolvedMath Contributors, UnsolvedMath: A Curated Collection of Open Mathematics Problems, catalogue entry. Catalogue curation and metadata: CC BY 4.0. The source extract is retained; definitions, English summaries and qualifications are editorial additions.
\url{https://huggingface.co/datasets/ulamai/UnsolvedMath}.
\item[S2] AIM Workshop Problem Lists, source linked by the catalogue.
\url{https://aimath.org/WWN/compalgstat/compalgstat.pdf}.
\end{itemize}
\paragraph{Preparation and review.}
The original source and UnsolvedMath contributors retain attribution for the source extract. Codex (GPT-6; exact serving revision not exposed) assisted with definitions, English summaries, scope notes and formatting. No new human mathematical authorship or independently verified solution is asserted. Catalogue identity checked against the complete supplied registry and current repository. Distinct upstream records, including repeated questions and broad or damaged extracts, are retained for complete coverage of the source CS category. A repeated record is not a new independent conjecture. Complete mathematical-equivalence and current-literature review has not been performed. Preparation and final review follow contribution prompts C and D, with source gaps disclosed. Structural and rendering checks do not establish a new theorem or open status.
\label{end:statement}

\section{Literature and context check}
AIM Question 49 is an extension agenda and fixes neither a ring nor a statistical model. The following two elementary models test concrete obstructions to transporting polynomial calculations unchanged. They are selected examples, not a resolution of the agenda.
\section{Research attempt: Denominator domains and positivity in Fourier coordinates}
First consider the rational equation $y/x=0$ on its specified domain $x\ne0$ in $\mathbb C^2$. Its solutions are $(x,0)$ with $x\ne0$. Clearing denominators gives $y=0$, which adds $(0,0)$ unless the domain condition is kept. More generally, on $x\ne0$ the equation $xy=0$ is equivalent to $y=0$, while in the polynomial ring the first zero set includes the entire line $x=0$. The saturation calculation here is exact:
\[
(xy):x^\infty=(y).
\]
Indeed $x y\in(xy)$ proves one inclusion. Conversely if $x^k f\in(xy)$, substitution $y=0$ gives $x^k f(x,0)=0$ in the integral domain $\mathbb Q[x]$, so $f(x,0)=0$. Expansion in powers of $y$ then shows $y\mid f$. Even saturation describes the closure, and the condition $x\ne0$ is still needed to describe the original rational domain.

For a Fourier example use the four-point group $G=\{-1,1\}^2$ with coordinatewise multiplication. For a probability mass function define $a=E S_1$, $b=E S_2$, $c=E S_1S_2$. The inverse formula is
\[
p(s_1,s_2)=\tfrac14(1+a s_1+b s_2+c s_1s_2).
\]
To verify it, sum the right-hand side against each of $1,s_1,s_2,s_1s_2$; every nonconstant sign character sums to zero and each character squared sums to four. Those four moment equations determine four masses because the character matrix times its transpose is $4I$. Thus the precise probability region in these coordinates is the four inequalities
\[
1+a s_1+b s_2+c s_1s_2\ge0\quad(s_1,s_2\in\{-1,1\}).
\]
The separate necessary bounds $|a|,|b|,|c|\le1$ are insufficient: $(a,b,c)=(9/10,9/10,-9/10)$ gives $p(-1,-1)=-17/40$.

There is nevertheless a genuine algebraic simplification. For independent group-valued $S,T$, each character satisfies $\chi(ST)=\chi(S)\chi(T)$, hence $E\chi(ST)=E\chi(S)E\chi(T)$. Convolution therefore becomes coordinatewise multiplication of the three Fourier moments. The failed extension route is to use that simple multiplication law and discard the coupled positivity inequalities; the displayed triple falsifies it.
\section{Verification and adversarial audit}
The Fourier masses sum to one for every triple; the counterexample fails only positivity. Its negative mass is $(1-9/10-9/10-9/10)/4=-17/40$. The rational example distinguishes localization, closure, and actual points in the domain.
\paragraph{Audit corrections.} The saturation is not equated with the punctured solution set. Fourier coordinate bounds are tested jointly rather than treated as sufficient probability constraints.
\paragraph{External computation.} None; symbolic arguments and boundary checks only.
\section{Scope and remaining gaps}
No ring or statistical target is fixed in the source. The calculations show exact domain and positivity conditions for two models; they do not choose a universally preferable ring or solve continuous Fourier or inference problems.
\paragraph{Integration review.} Codex (GPT-6; exact serving revision not exposed) reviewed the written arguments and expanded layout and source-reference wrappers. Repository metadata and the production web renderer were checked using Python 3.13 and JavaScript with MathJax 3.2.2. This model-assisted review is not independent human proof verification; no PDF compilation or formal proof checking is asserted.

\section{Completion Estimate}
\noindent\textbf{COMPLETION: 0\%}
The calculations below the preserved source statement investigate a concrete route,
but establish no substantiated new progress toward the entire original target.
Known special cases and failed sufficient conditions are not credited as a solution.
\section{References}
The original source references and catalogue attribution are preserved in the
Sources section above. The following URLs were actually consulted for this attempt;
retention of an additional original reference does not claim fresh consultation of it.
\begin{itemize}
\item \url{https://aimath.org/WWN/compalgstat/compalgstat.pdf} (consulted 2026-10-11). 
\end{itemize}
\section{Model and contributor disclosure}
OpenAI gpt-6-astra, requested reasoning configuration ultra, generated this research
attempt on 2026-10-11. The exact serving revision is not exposed. Codex GPT-6
performed orchestration, repository-format review, and a model-assisted mathematical
audit of the written arguments; its exact serving revision
is also not exposed. The mathematical route, boundary checks and final adversarial
self-audit followed contribution prompts A then B in
\texttt{docs/contribution\_prompts.md}, after reading \texttt{CONTRIBUTING.md}.
Python 3.14.2 assembled the standalone documents; any mathematical computation is
specified in Verification. No theorem prover was used. No human mathematical
authorship, substantial human proof correction, or independent human proof checking
is asserted. Model-assisted review, source retrieval, and structural checks do not
establish correctness or novelty. Full private reasoning transcripts are unavailable
as reproducibility artifacts; the complete arguments, inputs and outcomes actually
used are stated here or in the companion batch audit.
\section{Final status}
Unresolved as a new research contribution. This is a concrete route investigation,
not a claim of new complete resolution or a current-status certification. A known
answer reproduced for context does not become a new contribution. The catalogue's
source-status label is preserved; this attempt status does not assert that a
source-labelled solved question is open.
\end{document}
